\scanpage{113}
Soit $(\alpha_n)$ une suite \uline{décroissante} \inserted{et convergeant vers zéro} de nombres
positifs. Soit $\rho$ la borne inférieure des nombres
\deleted{positifs} $s>0$ tels que $(\alpha_n)$ soit majorée par
un multiple de la série $\left(\frac{1}{n^{1/s}}\right)$ (et $+\infty$
si un tel $s$ n'existe pas) et soit $\sigma$ la
borne inférieure des nombres $s>0$, tels que
$\sum_n\alpha_n^s<+\infty$. (et $+\infty$ si un tel $s$ n'existe pas)

On a aussi $\frac1\rho\deleted{\unsure{\leq}}\inserted{=}\limsup_n\unsure{:}\frac{\log(1/\alpha_n)}{\log n}$. \deleted{$\rho$ et $\sigma$
sont \unsure{resp\textsuperscript{t}} les bornes inf. d'ensembles
A et B.} -- Le critère de Riemann affirme
que \deleted{$\unsure{\rho\leq|\rho\leq\sigma}$} $\sigma\leq\rho$. Mais on a aussi $\rho\leq\sigma$,
car \unsure{on a le|en effet}

\uline{Lemme} \quad Si $s>0$ est tel que $\sum_n\alpha_n^s<+\infty$
alors \inserted{$\geq$} \deleted{la suite $(\alpha_n)$ est majorée par un
multiple de la suite $\left(\frac1{n^{1/s}}\right)$} \inserted{$\alpha_n=o\left(\frac1{n^{1/s}}\right)$}. ($\alpha_n$ étant décroissant !)

Démonstration. On se ramène au cas où $s=1$
(en posant $B_n=\alpha_n^s$) \inserted{et à prouver qu'il est impossible que} \deleted{\unsure{sup}} pour une suite
strictement croissante d'entiers $n_i$, on \deleted{avait} ait
$\alpha_{n_i}\geq\frac1{n_i}$ ; \inserted{en effet, alors} on aurait \quad $\alpha_{n_{i-1}+1}+\alpha_{n_{i-1}+2}\unsure{+|\mathord{-}}\,\mathord{-}+\alpha_{\unsure{n_i|n}}\geq\frac{n_i-n_{i-1}}{n_i}$
soit $\varepsilon_i$ ce second membre, $\sum\varepsilon_i$ est \unsure{donc}
convergente, et on a $n_i(1-\varepsilon_i)=n_{i-1}$, d'où $n_i=\frac{n_{i-1}}{1-\varepsilon_i}$
$\displaystyle =\frac{n_{i-2}}{(1-\varepsilon_i)(1-\varepsilon_{i-1})}\mathord{-}\,\mathord{-}\,\mathord{-}\,\mathord{-}=\frac1{(1-\varepsilon_{\unsure{1|i|0}})\unsure{\mathord{.}\,\mathord{-}|\ldots}(1-\varepsilon_{\unsure{n|i|m}})}\leq\frac1{\prod_1^\infty(1-\varepsilon_k)}$, la suite
des $n_i$ est donc majorée, et s'arrête donc.
\margin{114}
